% !TEX program = xelatex
\documentclass[11pt,a4paper]{article}
\usepackage[margin=2.2cm,top=2cm]{geometry}
\usepackage{fontspec}
\usepackage{polyglossia}
\setdefaultlanguage{turkish}
\usepackage{amsmath,amssymb}
\usepackage{graphicx}
\usepackage{booktabs}
\usepackage{enumitem}
\usepackage{multicol}
\usepackage[most]{tcolorbox}
\usepackage{fancyhdr}
\usepackage{titlesec}
\usepackage{xcolor}
\usepackage{needspace}
\usepackage[hidelinks]{hyperref}

\definecolor{anamavi}{RGB}{0,45,98}
\definecolor{cozumyesil}{RGB}{0,110,50}
\definecolor{uyarikirmizi}{RGB}{180,20,20}

\titleformat{\section}{\large\bfseries\color{anamavi}}{\thesection.}{0.5em}{}
\titleformat{\subsection}{\normalsize\bfseries\color{anamavi}}{\thesubsection}{0.5em}{}

\pagestyle{fancy}\fancyhf{}
\lhead{\small MMAT105 Genel Matematik}
\rhead{\small Hafta 13 -- Ders Notu}
\cfoot{\small\thepage}
\renewcommand{\headrulewidth}{0.4pt}

\newtcolorbox{kural}[1][Kural]{colback=anamavi!5,colframe=anamavi,title=#1,fonttitle=\bfseries,breakable,left=4pt,right=4pt,top=2pt,bottom=2pt}
\newtcolorbox{dikkat}[1][Dikkat]{colback=uyarikirmizi!5,colframe=uyarikirmizi,title=#1,fonttitle=\bfseries,breakable,left=4pt,right=4pt,top=2pt,bottom=2pt}
\newcounter{ornek}
\newenvironment{ornek}{\refstepcounter{ornek}\par\medskip\needspace{8\baselineskip}\noindent\textbf{\color{anamavi}Örnek \theornek.}\ }{\par}
\newenvironment{cozum}{\begin{tcolorbox}[colback=cozumyesil!6,colframe=cozumyesil,boxrule=0.5pt,left=4pt,right=4pt,top=2pt,bottom=2pt,breakable]\textbf{\color{cozumyesil}Çözüm.}\ }{\end{tcolorbox}}
\newcounter{alis}
\newenvironment{alistirma}{\refstepcounter{alis}\par\medskip\needspace{9\baselineskip}\noindent\textbf{\color{anamavi}Alıştırma \thealis.}\ }{\par}
\newcommand{\kendinizoncesi}{\needspace{12\baselineskip}}
\newcommand{\ab}[1]{\overline{#1}}
\setlength{\parindent}{0pt}
\setlength{\parskip}{3pt}

\begin{document}

% ---------------- Başlık ----------------
\begin{minipage}[c]{\textwidth}
{\LARGE\bfseries\color{anamavi} MMAT105 Genel Matematik}\\[4pt]
{\large Hafta 13 Ders Notu: Sayı ve Kesir Problemleri,\\ Yaş Problemleri}\\[6pt]
{\small 2026--2027 Güz}
\end{minipage}

\vspace{4pt}\noindent\textcolor{anamavi}{\rule{\textwidth}{1pt}}

\textbf{Bu notun kullanımı.} Her bölümde önce kısa bir kural kutusu, ardından çözümlü örnekler yer alır. Örnekleri çözümüne bakmadan önce kendiniz deneyin. Her çözümün sonunda bir \textbf{sağlama} vardır: bulduğunuz değeri soru metnine yerleştirip doğrulamayı alışkanlık haline getirin. Notun sonunda önce çözümlü alıştırmalar, ardından cevap anahtarlı ``Kendiniz Çözün'' alıştırmaları yer alır.

% =====================================================
\section{Sözel İfadeyi Denkleme Çevirme}

\begin{kural}[Çeviri tablosu (bilinmeyen sayı \(x\))]
\begin{center}
\renewcommand{\arraystretch}{1.35}
\begin{tabular}{@{}ll@{\qquad}ll@{}}
Bir sayının \(3\) katının \(5\) eksiği & \(3x-5\) & Bir sayının \(5\) eksiğinin \(3\) katı & \(3(x-5)\)\\
Bir sayının \(2\) katından \(7\) fazlası & \(2x+7\) & Bir sayının \(4\) fazlasının yarısı & \(\tfrac{x+4}{2}\)\\
Yarısı ile üçte birinin toplamı & \(\tfrac{x}{2}+\tfrac{x}{3}\) & Karesinin \(2\) katı & \(2x^2\)\\
\(a\), \(b\)'nin \(3\) katından \(2\) eksik & \(a=3b-2\) & \(a\) ile \(b\)'nin farkının karesi & \((a-b)^2\)
\end{tabular}
\end{center}
\textbf{Adımlar:} (1) bilinmeyeni adlandır, (2) diğer büyüklükleri onun cinsinden yaz, (3) denklemi kur ve çöz, (4) sağlama yap ve cevabın anlamlı olduğunu kontrol et.
\end{kural}

\begin{dikkat}
``\(5\) eksiğinin \(3\) katı'' ile ``\(3\) katının \(5\) eksiği'' farklıdır: önce yapılan işlem paranteze girer. Bulunan \(x\) her zaman istenen cevap değildir (``büyük sayı'', ``mevcut'' \dots).
\end{dikkat}

\begin{ornek}
Hangi sayının \(3\) katının \(5\) eksiği, aynı sayının yarısının \(10\) fazlasına eşittir?
\end{ornek}
\begin{cozum}
\[3x-5=\frac{x}{2}+10\ \xrightarrow{\ \cdot 2\ }\ 6x-10=x+20\ \Rightarrow\ 5x=30\ \Rightarrow\ x=6\]
Sağlama: \(3\cdot 6-5=13\), \ \(\frac{6}{2}+10=13\) \checkmark
\end{cozum}

\begin{ornek}
Bir sayının \(2\) fazlasının \(3\) katı, aynı sayının \(3\) katının \(2\) fazlasından kaç fazladır?
\end{ornek}
\begin{cozum}
\(3(x+2)-(3x+2)=3x+6-3x-2=4\). Sonuç \(x\)'e bağlı değildir: her sayı için fark \(4\)'tür.\\
(Örneğin \(x=10\): \(36-32=4\).) Parantezin önemi bu örnekte açıkça görülür.
\end{cozum}

\begin{ornek}
Bir otoparkta yalnızca otomobil (\(4\) tekerlekli) ve motosiklet (\(2\) tekerlekli) vardır. Toplam \(30\) araç ve \(100\) tekerlek olduğuna göre kaç motosiklet vardır?
\end{ornek}
\begin{cozum}
Otomobil \(o\), motosiklet \(m\): \quad \(o+m=30\), \ \(4o+2m=100\). \ İkinci denklemi \(2\)'ye bölüp birinciyi çıkaralım:
\[2o+m=50,\ \ o+m=30\ \Rightarrow\ o=20,\ m=10\]
Sağlama: \(80+20=100\) tekerlek \checkmark \quad Cevap: \(10\) motosiklet.
\end{cozum}

\begin{ornek}
Bir sınıftaki öğrenciler sıralara \(2\)'şerli oturursa \(6\) öğrenci ayakta kalıyor; \(3\)'erli oturursa \(2\) sıra boş kalıyor (diğer sıralar tam dolu). Sınıf mevcudu kaçtır?
\end{ornek}
\begin{cozum}
Sıra sayısı \(x\): \quad \(2x+6=3(x-2)\ \Rightarrow\ 2x+6=3x-6\ \Rightarrow\ x=12\).\\
Mevcut \(=2\cdot 12+6=30\).\\
Sağlama: \(3\)'erli oturuşta \(10\) sıra dolu: \(3\cdot 10=30\) \checkmark
\end{cozum}

\begin{ornek}
Bir merdivenin basamaklarını \(3\)'er \(3\)'er çıkıp \(2\)'şer \(2\)'şer inen bir kişinin inerken attığı adım sayısı, çıkarken attığı adım sayısından \(6\) fazladır. Merdiven kaç basamaklıdır?
\end{ornek}
\begin{cozum}
Basamak sayısı hem \(2\)'ye hem \(3\)'e bölünmeli: \(6k\) diyelim. Çıkış: \(\frac{6k}{3}=2k\) adım; iniş: \(\frac{6k}{2}=3k\) adım.
\[3k-2k=6\ \Rightarrow\ k=6\ \Rightarrow\ 36\ \text{basamak}\]
Sağlama: \(18-12=6\) \checkmark
\end{cozum}

% =====================================================
\section{Ardışık Sayı ve Basamak Problemleri}

\begin{kural}[Ardışık sayılar ve basamak çözümlemesi]
Ardışık tam sayılar: \(x,\ x+1,\ x+2,\dots\) \quad Ardışık çift/tek sayılar: \(x,\ x+2,\ x+4,\dots\)\\
Tek sayıda ardışık terimin toplamı \(=\) terim sayısı \(\times\) ortanca terim.\\[2pt]
\(a,b\) rakam ve \(a\neq 0\) olmak üzere: \ \(\ab{ab}=10a+b\), \quad \(\ab{ab}+\ab{ba}=11(a+b)\), \quad \(\ab{ab}-\ab{ba}=9(a-b)\).\\
Bir \(N\) sayısının sağına \(c\) rakamı yazılırsa yeni sayı \(10N+c\) olur.
\end{kural}

\begin{dikkat}
\(\ab{ab}\) bir çarpım değildir: \(\ab{ab}\neq a\cdot b\). Rakamlar \(0\)--\(9\) arasındadır; bulduğunuz \(a\), \(b\) değerlerinin rakam olduğunu kontrol edin.
\end{dikkat}

\begin{ornek}
Ardışık beş tam sayının toplamı \(130\)'dur. En büyük ile en küçük sayının çarpımı kaçtır?
\end{ornek}
\begin{cozum}
Ortanca terim \(130\div 5=26\) \(\Rightarrow\) sayılar \(24,25,26,27,28\). \quad \(24\cdot 28=672\).
\end{cozum}

\begin{ornek}
İki basamaklı \(\ab{ab}\) sayısının rakamları toplamı \(12\)'dir. Bu sayı, rakamlarının yerleri değiştirilerek elde edilen sayıdan \(36\) fazladır. \(\ab{ab}\) kaçtır?
\end{ornek}
\begin{cozum}
\(\ab{ab}-\ab{ba}=9(a-b)=36\Rightarrow a-b=4\). \quad \(a+b=12\) ile: \(a=8,\ b=4\) \(\Rightarrow\) \(\ab{ab}=84\).\\
Sağlama: \(84-48=36\) \checkmark
\end{cozum}

\begin{ornek}
İki basamaklı bir sayının sağına \(3\) rakamı yazılınca elde edilen üç basamaklı sayı, ilk sayıdan \(444\) fazla oluyor. İlk sayı kaçtır?
\end{ornek}
\begin{cozum}
Sayı \(N\) ise yeni sayı \(10N+3\):
\[10N+3-N=444\ \Rightarrow\ 9N=441\ \Rightarrow\ N=49\]
Sağlama: \(493-49=444\) \checkmark
\end{cozum}

% =====================================================
\section{Kesir Problemleri}

\begin{kural}[Bütünün kesri ve kalanın kesri; depo doluluk oranı]
Bir \(A\) büyüklüğünün \(\frac{m}{n}\)'si \(\frac{m}{n}\cdot A\)'dır; \(\frac{m}{n}\)'si \(B\) ise \(A=B\cdot\frac{n}{m}\).\\
\textbf{Kalanın kesri:} önce kalan bulunur, kesir kalana uygulanır.\\
\textbf{EKOK yöntemi:} Bütüne \(x\) demek yerine paydaların EKOK'u kadar birim verin (\(12k\), \(20k\), \dots); her parça tam sayı olur.\\
\textbf{Depo:} doluluk \(\frac{a}{b}\)'den \(\frac{c}{d}\)'ye çıkıyorsa eklenen su \(\left(\frac{c}{d}-\frac{a}{b}\right)V\)'dir.
\end{kural}

\begin{dikkat}
``Maaşının \(\frac{1}{5}\)'ini, sonra \(\frac{1}{3}\)'ünü harcadı'' (iki kesir de maaşın) ile ``\(\frac{1}{5}\)'ini, sonra \textbf{kalanın} \(\frac{1}{3}\)'ünü harcadı'' farklı sorulardır. Soruyu dikkatle okuyun.
\end{dikkat}

\begin{ornek}
Bir paranın önce \(\frac{1}{4}\)'ü, sonra kalanın \(\frac{1}{3}\)'ü harcanıyor. Geriye \(100\) TL kaldığına göre paranın tamamı kaç TL'dir?
\end{ornek}
\begin{cozum}
Para \(=12k\). \ \(\frac{1}{4}\)'ü: \(3k\) \(\to\) kalan \(9k\); \ kalanın \(\frac{1}{3}\)'ü: \(3k\) \(\to\) kalan \(6k=100\).
\[k=\frac{50}{3},\qquad 12k=200\ \text{TL}\]
Sağlama: \(200\to 150\to 100\) \checkmark \ (\(k\)'nin tam sayı çıkması gerekmez; sonuç tam sayıdır.)
\end{cozum}

\begin{ornek}
Bir top bırakıldığı yüksekliğin \(\frac{3}{4}\)'ü kadar zıplıyor. Top yere ikinci vuruşundan sonra \(18\) cm yükseldiğine göre ilk bırakıldığı yükseklik kaç cm'dir?
\end{ornek}
\begin{cozum}
Sondan başa: \(18\cdot\frac{4}{3}=24\) cm (birinci zıplama), \ \(24\cdot\frac{4}{3}=32\) cm.\\
Sağlama: \(32\to 24\to 18\) \checkmark
\end{cozum}

\begin{ornek}
Bir deponun \(\frac{3}{8}\)'i doludur. Depoya \(45\) litre su eklenince deponun yarısı doluyor. Deponun tamamı kaç litre su alır?
\end{ornek}
\begin{cozum}
\(\frac{1}{2}-\frac{3}{8}=\frac{1}{8}\) \(\Rightarrow\) \(\frac{1}{8}V=45\Rightarrow V=360\) L.
\quad Sağlama: \(135+45=180=\frac{1}{2}\cdot 360\) \checkmark
\end{cozum}

\begin{ornek}
Bir telin bir ucundan \(\frac{1}{6}\)'sı kesilirse telin orta noktası \(5\) cm kayıyor. Telin kesilmeden önceki boyu kaç cm'dir?
\end{ornek}
\begin{cozum}
Tel \(L\) cm. Önce orta nokta kalan uçtan \(\frac{L}{2}\) uzakta; kesimden sonra boy \(\frac{5L}{6}\), orta nokta \(\frac{5L}{12}\) uzakta.
\[\frac{L}{2}-\frac{5L}{12}=\frac{L}{12}=5\ \Rightarrow\ L=60\ \text{cm}\]
Kısa yol: orta nokta, kesilen parçanın yarısı kadar kayar: \(\frac{1}{2}\cdot\frac{L}{6}=5\).
\end{cozum}

\begin{ornek}
Pay ile paydasının toplamı \(36\) olan bir kesrin payından \(2\) çıkarılıp paydasına \(2\) eklenince kesir \(\frac{1}{3}\)'e eşit oluyor. Bu kesir nedir?
\end{ornek}
\begin{cozum}
Pay \(p\), payda \(36-p\):
\[\frac{p-2}{38-p}=\frac{1}{3}\ \Rightarrow\ 3p-6=38-p\ \Rightarrow\ 4p=44\ \Rightarrow\ p=11\]
Kesir \(\frac{11}{25}\). \quad Sağlama: \(\frac{9}{27}=\frac{1}{3}\) \checkmark
\end{cozum}

% =====================================================
\section{Yaş Problemleri}

\begin{kural}[Temel kurallar ve tablo yöntemi]
\begin{itemize}[leftmargin=*,itemsep=1pt]
\item İki kişi arasındaki yaş farkı hiç değişmez; yaşlar oranı (kat) ise zamanla değişir.
\item \(t\) yıl sonra herkesin yaşı \(t\) artar; \(n\) kişinin yaşları toplamı \(n\cdot t\) artar, yaş ortalaması \(t\) artar.
\item ``\(A\), \(B\)'nin yaşındayken'' \(=\) \(A-B\) yıl önce; ``\(B\), \(A\)'nın yaşına geldiğinde'' \(=\) \(A-B\) yıl sonra.
\item Tablo: satırlara kişiler, sütunlara \textit{\(t\) yıl önce / Şimdi / \(t\) yıl sonra}; denklem, sorunun söylediği sütundan kurulur.
\item Doğum yılı soruları: yaş \(=\) (yıl \(-\) doğum yılı) (doğum günü geçmiş kabul edilir).
\end{itemize}
\end{kural}

\begin{dikkat}
Sonucu mutlaka kontrol edin: geçmişteki yaşlar negatif olamaz; bir annenin, çocuğundan küçük ya da çok yakın yaşta olduğu bir sonuç kurulan denklemde hata olduğunu gösterir.
\end{dikkat}

\begin{ornek}
Bir babanın yaşı, üç çocuğunun yaşları toplamının \(3\) katıdır. \(2\) yıl sonra babanın yaşı, çocukların yaşları toplamının \(2\) katı olacaktır. Baba bugün kaç yaşındadır?
\end{ornek}
\begin{cozum}
Çocukların yaş toplamı \(s\), baba \(3s\). \(2\) yıl sonra toplam \(3\cdot 2=6\) artar:
\[3s+2=2(s+6)\ \Rightarrow\ 3s+2=2s+12\ \Rightarrow\ s=10,\quad \text{baba}=30\]
Sağlama: \(32=2\cdot 16\) \checkmark
\end{cozum}

\begin{ornek}
\(3\) yıl önceki yaşlarının ortalaması \(18\) olan \(4\) kişilik bir ailenin \(2\) yıl sonraki yaşları toplamı kaç olur?
\end{ornek}
\begin{cozum}
\(3\) yıl önceden \(2\) yıl sonrasına \(5\) yıl geçer; ortalama \(5\) artar: \(23\). \quad Toplam \(=23\cdot 4=92\).\\
(Diğer yol: \(18\cdot 4=72\), \ \(72+4\cdot 5=92\).)
\end{cozum}

\begin{ornek}
Bir annenin yaşı, iki çocuğunun yaşları farkının \(5\) katıdır. \(8\) yıl sonra annenin yaşı, çocukların yaşları farkının \(6\) katı olacaktır. Anne bugün kaç yaşındadır?
\end{ornek}
\begin{cozum}
Çocukların yaş farkı \(d\) zamanla \textbf{değişmez}. Anne \(5d\); \(8\) yıl sonra \(5d+8=6d\Rightarrow d=8\). \quad Anne \(=40\).\\
Sağlama: \(8\) yıl sonra anne \(48=6\cdot 8\) \checkmark
\end{cozum}

\begin{ornek}
Bir dedenin yıl olarak yaşı, torununun ay olarak yaşına eşittir. Dede ile torununun yıl olarak yaşları toplamı \(65\) olduğuna göre torun kaç yaşındadır?
\end{ornek}
\begin{cozum}
Torun \(t\) yaşında ise ay olarak yaşı \(12t\); dede \(12t\) yaşında. \quad \(12t+t=65\Rightarrow t=5\), dede \(60\).
\end{cozum}

\begin{ornek}
Can, Ece'nin bugünkü yaşındayken Ece \(6\) yaşındaydı. Ece, Can'ın bugünkü yaşına geldiğinde Can \(30\) yaşında olacaktır. Can ile Ece bugün kaç yaşındadır?
\end{ornek}
\begin{cozum}
Yaş farkı \(d\) olsun. Sayı doğrusunda \(6,\ E,\ C,\ 30\) değerleri aralarında \(d\) olacak şekilde sıralanır:
\[6+3d=30\ \Rightarrow\ d=8\ \Rightarrow\ E=14,\ C=22\]
Sağlama: \(8\) yıl önce Can \(14\), Ece \(6\); \(8\) yıl sonra Ece \(22\), Can \(30\) \checkmark
\end{cozum}

\begin{ornek}
Bir anne 1978, kızı 2004 yılında doğmuştur. Hangi yılda annenin yaşı kızının yaşının \(3\) katı olur?
\end{ornek}
\begin{cozum}
Yaş farkı \(2004-1978=26\) sabittir. Kız \(k\) yaşındayken anne \(k+26\):
\[k+26=3k\ \Rightarrow\ k=13\ \Rightarrow\ 2004+13=2017\]
Sağlama: 2017'de anne \(39=3\cdot 13\) \checkmark
\end{cozum}

\begin{ornek}
\(5\) kişilik bir grubun yaş ortalaması \(24\)'tür. Gruba bir kişi katılıyor; \(3\) yıl sonra bu \(6\) kişinin yaş ortalaması \(28\) olacaktır. Gruba katılan kişi bugün kaç yaşındadır?
\end{ornek}
\begin{cozum}
\(6\) kişinin bugünkü ortalaması \(28-3=25\) \(\Rightarrow\) toplam \(25\cdot 6=150\). \ İlk \(5\) kişinin toplamı \(24\cdot 5=120\).
\[\text{Katılan kişi}=150-120=30\]
\end{cozum}

% =====================================================
\section{Özet: Formül Kartı}
\begin{center}
\renewcommand{\arraystretch}{1.6}
\begin{tabular}{@{}ll@{}}
\toprule
Konu & Formül / Kural\\\midrule
Çeviri & ``\(a\) katının \(b\) eksiği'' \(=ax-b\); \ ``\(b\) eksiğinin \(a\) katı'' \(=a(x-b)\)\\
Toplam ve fark & büyük \(=\dfrac{T+F}{2}\), \ küçük \(=\dfrac{T-F}{2}\)\\
Ardışık sayılar & \(x,\ x+1,\dots\); \ tek/çift: \(x,\ x+2,\dots\); \ toplam \(=n\cdot\)ortanca (\(n\) tek)\\
Basamak & \(\ab{ab}=10a+b\); \ \(\ab{ab}+\ab{ba}=11(a+b)\); \ \(\ab{ab}-\ab{ba}=9(a-b)\)\\
Sağına rakam yazma & \(N\to 10N+c\)\\
Kesir & \(\frac{m}{n}\)'si \(B\) ise bütün \(=B\cdot\frac{n}{m}\); \ bütün \(=\) paydaların EKOK'u kadar birim\\
Kalanın kesri & önce kalan, sonra kesir\\
Depo & \(\left(\frac{c}{d}-\frac{a}{b}\right)V=\) eklenen (kullanılan) miktar\\
Yaş farkı & hiç değişmez; oran değişir\\
Yaş toplamı & \(n\) kişi, \(t\) yıl \(\Rightarrow\) toplam \(+\,n\cdot t\), ortalama \(+\,t\)\\
``Senin yaşındayken'' & geçen süre \(=\) yaş farkı \(d\); \ \(B-d,\ B,\ A,\ A+d\) eşit aralıklı\\\bottomrule
\end{tabular}
\end{center}

% =====================================================
\newpage
% Hafta 13 -- alıştırma gövdesi (ders notu ve ayrı alıştırma kağıdı tarafından \input edilir)
\section{Çözümlü Alıştırmalar}
Her soruyu önce kendiniz deneyin, sonra çözümü okuyun.

\begin{alistirma}
Ardışık üç çift sayının toplamı \(138\)'dir. En küçüğü kaçtır?
\end{alistirma}
\begin{cozum}
Sayılar \(x,\ x+2,\ x+4\) (\(x\) çift):
\[x+(x+2)+(x+4)=138\ \Rightarrow\ 3x+6=138\ \Rightarrow\ x=44\]
Sayılar \(44,\ 46,\ 48\); toplam \(138\) \checkmark. \ (Kısa yol: ortanca \(138\div 3=46\).)
\end{cozum}

\begin{alistirma}
Toplamları \(85\) olan iki doğal sayıdan büyüğü küçüğüne bölündüğünde bölüm \(3\), kalan \(9\) oluyor. Büyük sayı kaçtır?
\end{alistirma}
\begin{cozum}
Küçük sayı \(k\) ise büyük sayı \(3k+9\) (bölme: bölünen \(=\) bölen \(\cdot\) bölüm \(+\) kalan).
\[k+3k+9=85\ \Rightarrow\ 4k=76\ \Rightarrow\ k=19,\qquad \text{büyük}=3\cdot 19+9=66\]
Sağlama: \(66=19\cdot 3+9\) ve kalan bölenden küçük (\(9<19\)) \checkmark
\end{cozum}

\begin{alistirma}
İki basamaklı bir sayının rakamları toplamı \(9\)'dur. Rakamlarının yerleri değiştirilince sayı \(45\) azalıyor. Bu sayı kaçtır?
\end{alistirma}
\begin{cozum}
\(\ab{ab}-\ab{ba}=(10a+b)-(10b+a)=9(a-b)=45\Rightarrow a-b=5\).
\[a+b=9,\ \ a-b=5\ \Rightarrow\ 2a=14\ \Rightarrow\ a=7,\ b=2\ \Rightarrow\ \ab{ab}=72\]
Sağlama: \(72-27=45\) \checkmark
\end{cozum}

\begin{alistirma}
Bir kişi parasının \(\frac{3}{8}\)'ini harcıyor, kalan parasının \(\frac{2}{5}\)'ini de kardeşine veriyor. Geriye \(180\) TL kaldığına göre başlangıçta kaç TL'si vardı?
\end{alistirma}
\begin{cozum}
Paydalar \(8\) ve \(5\) \(\Rightarrow\) para \(=40k\).\\
Harcanan: \(\frac{3}{8}\cdot 40k=15k\) \(\to\) kalan \(25k\); \ kardeşine: \(\frac{2}{5}\cdot 25k=10k\) \(\to\) kalan \(15k\).
\[15k=180\ \Rightarrow\ k=12\ \Rightarrow\ 40k=480\ \text{TL}\]
Sağlama: \(480-180=300\), \ \(300-120=180\) \checkmark
\end{cozum}

\begin{alistirma}
Bir su deposunun \(\frac{1}{3}\)'ü doludur. Depoya \(50\) litre su eklenince deponun \(\frac{3}{4}\)'ü doluyor. Deponun tamamı kaç litre su alır?
\end{alistirma}
\begin{cozum}
Eklenen su, doluluk oranlarının farkına karşılık gelir:
\[\frac{3}{4}-\frac{1}{3}=\frac{9-4}{12}=\frac{5}{12}\ \Rightarrow\ \frac{5}{12}V=50\ \Rightarrow\ V=50\cdot\frac{12}{5}=120\ \text{L}\]
Sağlama: \(\frac{1}{3}\cdot 120=40\), \ \(40+50=90=\frac{3}{4}\cdot 120\) \checkmark
\end{cozum}

\begin{alistirma}
İki kardeşin yaşları toplamı \(30\)'dur. \(5\) yıl önce büyük kardeşin yaşı küçüğün yaşının \(3\) katıydı. Kardeşler bugün kaç yaşındadır?
\end{alistirma}
\begin{cozum}
Tablo yöntemi (küçüğün yaşı \(k\), büyüğünki \(30-k\)):
\begin{center}
\begin{tabular}{lcc}
 & \(5\) yıl önce & Şimdi\\\hline
Küçük & \(k-5\) & \(k\)\\
Büyük & \(25-k\) & \(30-k\)
\end{tabular}
\end{center}
\[25-k=3(k-5)\ \Rightarrow\ 25-k=3k-15\ \Rightarrow\ 4k=40\ \Rightarrow\ k=10\]
Küçük \(10\), büyük \(20\) yaşında. \quad Sağlama: \(5\) yıl önce \(15=3\cdot 5\) \checkmark
\end{cozum}

\begin{alistirma}
Bir baba \(46\) yaşında, üç çocuğu \(14\), \(10\) ve \(6\) yaşındadır. Kaç yıl sonra babanın yaşı çocuklarının yaşları toplamına eşit olur?
\end{alistirma}
\begin{cozum}
Çocukların yaş toplamı bugün \(30\); her yıl \(3\) artar (\(3\) kişi). Babanın yaşı her yıl \(1\) artar.
\[46+t=30+3t\ \Rightarrow\ 2t=16\ \Rightarrow\ t=8\]
Sağlama: \(8\) yıl sonra baba \(54\); çocuklar \(22+18+14=54\) \checkmark
\end{cozum}

\begin{alistirma}
Mert \(32\) yaşındadır. Mert, Ozan'ın bugünkü yaşındayken Ozan \(20\) yaşındaydı. Ozan, Mert'in bugünkü yaşına geldiğinde Mert kaç yaşında olur?
\end{alistirma}
\begin{cozum}
Ozan \(x\), yaş farkı \(d=32-x\). Mert'in \(x\) yaşında olduğu an \(d\) yıl öncedir; o zaman Ozan \(x-d=20\):
\[x-(32-x)=20\ \Rightarrow\ 2x=52\ \Rightarrow\ x=26,\quad d=6\]
Ozan'ın \(32\) yaşına gelmesi \(6\) yıl sürer; o zaman Mert \(32+6=38\) yaşında olur.\\
Sağlama: \(20,\ 26,\ 32,\ 38\) eşit aralıklı (\(d=6\)) \checkmark
\end{cozum}

\kendinizoncesi
\section{Alıştırmalar (Kendiniz Çözün)}

\begin{enumerate}[label=\textbf{\arabic*.},leftmargin=*,itemsep=6pt]
\item Bir sayının \(4\) katının \(7\) fazlası \(43\)'tür. Bu sayı kaçtır?
\item Bir sayının \(3\) eksiğinin \(5\) katı, aynı sayının \(2\) katının \(6\) fazlasına eşittir. Bu sayı kaçtır?
\item İki sayının toplamı \(75\), farkı \(21\)'dir. Küçük sayı kaçtır?
\item Ardışık beş tek sayının toplamı \(115\)'tir. En büyüğü kaçtır?
\item \(\ab{ab}\) ve \(\ab{ba}\) iki basamaklı sayılardır. \(\ab{ab}+\ab{ba}=154\) ve \(a-b=2\) ise \(\ab{ab}\) kaçtır?
\item Bir sınıftaki öğrenciler sıralara \(4\)'erli oturursa \(2\) sıra boş kalıyor (diğer sıralar tam dolu); \(3\)'erli oturursa \(5\) öğrenci ayakta kalıyor. Sınıf mevcudu kaçtır?
\item Bir sayının \(\frac{3}{5}\)'i \(27\) ise bu sayının \(\frac{2}{9}\)'u kaçtır?
\item Bir kişi parasının \(\frac{1}{6}\)'sını harcıyor, kalan parasının \(\frac{2}{5}\)'ini de arkadaşına veriyor. Geriye \(90\) TL kaldığına göre başlangıçta kaç TL'si vardı?
\item Bir deponun \(\frac{5}{8}\)'i doludur. Depodan \(36\) litre su kullanılınca deponun \(\frac{1}{4}\)'ü dolu kalıyor. Deponun tamamı kaç litredir?
\item Bir sayının \(\frac{3}{4}\)'ünün \(2\) eksiği, aynı sayının \(\frac{2}{3}\)'ünün \(1\) fazlasına eşittir. Bu sayı kaçtır?
\item Bir anne \(34\), oğlu \(7\) yaşındadır. Kaç yıl sonra annenin yaşı oğlunun yaşının \(4\) katı olur?
\item \(6\) kişilik bir ailenin yaş ortalaması \(20\)'dir. Ailenin en küçük üyesi \(7\) yaşında olduğuna göre \(5\) yıl önce bu ailenin yaşları toplamı kaçtı?
\item Bir babanın yaşı kızının yaşının \(5\) katıdır. \(3\) yıl sonra babanın yaşı kızının yaşının \(4\) katı olacaktır. Baba bugün kaç yaşındadır?
\item Can, Ece'den \(6\) yaş büyüktür. Ece, Can'ın bugünkü yaşına geldiğinde ikisinin yaşları toplamı \(50\) olacaktır. Ece bugün kaç yaşındadır?
\item Bir baba 1992'de, oğlu 2019'da doğmuştur. Hangi yılda babanın yaşı oğlunun yaşının \(4\) katı olur?
\end{enumerate}

\subsection*{Cevap Anahtarı}
{\small
\begin{multicols}{2}
\begin{enumerate}[label=\textbf{\arabic*.},leftmargin=*,itemsep=1pt]
\item \(4x+7=43\Rightarrow x=9\).
\item \(5x-15=2x+6\Rightarrow x=7\).
\item \((75-21)/2=27\) \ (büyük \(48\)).
\item Ortanca \(115/5=23\); sayılar \(19,\dots,27\) \(\Rightarrow\) \(27\).
\item \(11(a+b)=154\Rightarrow a+b=14\); \(a=8,b=6\) \(\Rightarrow\) \(86\).
\item \(4(x-2)=3x+5\Rightarrow x=13\) sıra; mevcut \(44\).
\item Sayı \(27\cdot\frac{5}{3}=45\); \(\frac{2}{9}\cdot 45=10\).
\item Para \(30k\): \(5k\), kalan \(25k\); \(10k\), kalan \(15k=90\Rightarrow 180\) TL.
\item \(\frac{5}{8}-\frac{1}{4}=\frac{3}{8}\); \(\frac{3}{8}V=36\Rightarrow V=96\) L.
\item \(\cdot 12\): \(9x-24=8x+12\Rightarrow x=36\).
\item \(34+t=4(7+t)\Rightarrow t=2\) \ (\(36=4\cdot 9\)).
\item Bugün \(120\); \(120-6\cdot 5=90\).
\item \(5x+3=4(x+3)\Rightarrow x=9\); baba \(45\).
\item \(6\) yıl sonra: Ece \(E+6\), Can \(E+12\); \(2E+18=50\Rightarrow E=16\).
\item Fark \(27\); \(4k=k+27\Rightarrow k=9\Rightarrow 2028\).
\end{enumerate}
\end{multicols}}


\vspace{6pt}
\noindent\textcolor{anamavi}{\rule{\textwidth}{0.5pt}}\\
{\small Gelecek hafta: Yüzde, kâr--zarar, karışım ve faiz problemleri; genel tekrar.}

\end{document}
